% ECONOMICS_PROBLEM: {"id": "F36-407", "title": "International risk sharing", "jel_code": "F36", "source_id": "OP-0335"}
% FIRST_POSTED: 2026-10-09
% ATTEMPT_STATUS: partial
% ENTRY_KIND: research_attempt
\documentclass[11pt]{article}
\usepackage[margin=1in]{geometry}
\usepackage[utf8]{inputenc}
\usepackage{amsmath,amssymb,amsthm,hyperref}
\newtheorem{theorem}{Theorem}
\newtheorem{proposition}[theorem]{Proposition}
\newtheorem{lemma}[theorem]{Lemma}
\newtheorem{corollary}[theorem]{Corollary}
\theoremstyle{definition}
\newtheorem{example}[theorem]{Example}
\newtheorem{remark}[theorem]{Remark}
\newcommand{\E}{\mathbb E}
\newcommand{\R}{\mathbb R}
\title{International risk sharing}
\author{GPT 6 Astra Ultra}
\date{9 October 2026}
\begin{document}
\maketitle
\section*{Target and the measurement problem}
The target is the change in a country-specific residual variance of log
consumption growth under an integration intervention. This is not itself a
welfare measure. We establish an exact residual-variance decomposition, solve
a finite-state three-country insurance benchmark, and identify what a policy
experiment would need to observe. The monetary, safe-asset and nominal-rigidity
extension remains open in this attempt.

\section*{Fixing the projection before comparing policies}
Let $g(a)$ be log consumption growth in regime $a$, and let $Z$ be a vector
of common predictors with an intercept removed after demeaning. In each regime
assume second moments and nonsingular $\Sigma_{ZZ}(a)$. The population linear
projection has coefficient $b(a)=\Sigma_{ZZ}(a)^{-1}\Sigma_{Zg}(a)$.
Its residual variance is
\[
 v(a)=\operatorname{Var}(g(a))-
 \Sigma_{gZ}(a)\Sigma_{ZZ}(a)^{-1}\Sigma_{Zg}(a).           \tag{1}
\]
\begin{proposition}
If instead one holds a baseline coefficient $b_0$ fixed, the residual variance
in regime $a$ equals
\[
 v(a)+(b_0-b(a))^{\top}\Sigma_{ZZ}(a)(b_0-b(a)).      \tag{2}
\]
\end{proposition}
\begin{proof}
Subtract the regime mean and write
$g-b_0^{\top}Z=(g-b(a)^{\top}Z)+(b(a)-b_0)^{\top}Z$.
The first summand is orthogonal to every coordinate of $Z$ by its normal
equations. Taking variances proves (2), and minimizing the quadratic expression
in $b$ gives (1).
\end{proof}
Thus a policy that changes exposure to a common shock can appear to worsen
insurance simply because an obsolete predictor is held fixed. Either convention
is definable, but they estimate different quantities. At singular covariance,
use a fixed predictor subspace and orthogonal projection or a stated generalized
inverse; arbitrary regularization changes the target unless its error is controlled.

A second elementary trap is welfare. Let $g(0)=m_0+e_0$ and
$g(1)=m_1+e_1$ with lower residual variance in regime one but $m_1<m_0$.
For one-period log utility, expected utility is determined by the mean of log
consumption, so welfare falls regardless of the lower variance. In a concrete
positive-consumption example, baseline $c$ is $1$ or $2$ equally often and
integrated $c$ is constantly $0.5$. Residual log variance falls to zero while
expected log consumption falls from $\tfrac12\log2$ to $-\log2$.
This is a logical example, not a feasible insurance policy with unchanged
aggregate resources; it shows why resource and mean effects must be separately
modeled before interpreting the statistic.

\section*{An exact three-country complete-market benchmark}
Label countries $U,A,E$. There is one representative household per country,
one consumption good, a baseline date with consumption one, and eight equally
likely date-one states $s=(s_U,s_A,s_E)\in\{-1,1\}^3$. Endowments are
$y_i(s)=1+h s_i$ with $0<h<1$. Households have identical log utility. Compare
autarky with complete contingent-claim trade, no transaction costs and common
beliefs. Identical initial expected wealth under symmetric equilibrium prices
supports equal Pareto weights. The state resource constraint is
$\sum_i c_i(s)=3+h\sum_i s_i$.

\begin{theorem}
The equal-weight efficient allocation, supportable by complete contingent
claims, is
\[
 c_i^*(s)=1+\frac h3(s_U+s_A+s_E)>0.                     \tag{3}
\]
Every country has weakly greater ex ante log utility than under autarky,
strictly greater when $h>0$.
\end{theorem}
\begin{proof}
For each state, strict concavity and the first-order conditions
$1/c_i=\lambda_s$ imply equal consumption, giving (3). Positive Arrow prices
proportional to state probability divided by $c_i^*(s)$ support individual
optimality. Symmetry of the endowment law and prices implies each country has
one third of aggregate state-priced wealth, so these consumptions satisfy its
budget. Statewise Jensen gives
$3\log(\sum_i y_i/3)\ge\sum_i\log y_i$, with strict inequality whenever
endowments differ. Take expectations and use exchangeability to conclude the
same ex ante improvement for each country. Mixed-sign states have positive
probability for $h>0$, making the inequality strict.
\end{proof}
With only a constant predictor, the autarky variance of log growth is
\[
 v_0=\frac14\left[\log\frac{1+h}{1-h}\right]^2.
\]
Under integration, log consumption takes values
$\log(1-h),\log(1-h/3),\log(1+h/3),\log(1+h)$ with probabilities
$1/8,3/8,3/8,1/8$. Therefore its exact residual variance is
\[
 v_1=\sum_{j=1}^4 p_j\ell_j^2-
                  \left(\sum_{j=1}^4p_j\ell_j\right)^2.  \tag{4}
\]
These finite sums determine the catalogue statistic without linearizing log
consumption. If the declared common predictor space contains all functions of
aggregate endowment, the integrated residual is instead identically zero,
because (3) is a function of that aggregate. The choice of information space
is consequently economically material even in this simple solved benchmark.

\section*{Asymmetric risk bearing and its limits}
For a level-consumption complete-market benchmark with CARA utility
$-\exp(-\gamma_i c_i)$ and positive Pareto weights $\lambda_i$, the same
statewise calculation, at an interior allocation, yields
\[
 c_i(s)=a_i+\theta_iY(s),\quad
 \theta_i=\frac{\gamma_i^{-1}}{\sum_j\gamma_j^{-1}},\quad
 \sum_i a_i=0,\qquad Y(s)=\sum_i y_i(s).                 \tag{5}
\]
To verify, the first-order conditions are
$\lambda_i\gamma_i\exp(-\gamma_i c_i)=\kappa_s$.
Take logarithms, divide by $\gamma_i$, sum over countries and solve for
$\log\kappa_s$; the constants produce $a_i$ and the coefficient of $Y$
produces $\theta_i$. On a finite state space choose endowments and weights
so all $c_i$ are positive; otherwise nonnegativity constraints create corners
and (5) is not valid. More risk-tolerant countries bear more aggregate level
risk, but (5) is not a formula for log residual variance or for safe-dollar
bond demand. It explicitly assumes away the incomplete markets and nominal
features of the source-motivated extension.

\section*{Identification and next empirical step}
Random assignment of an integration rule across independent three-country
clusters would identify the law of $(g_i,Z)$ in each regime and therefore (1),
provided exposure remains inside clusters, consumption measurement is comparable,
and second moments exist. A single worldwide policy episode generally provides
no such independent experiment. Instrumental exposure designs additionally need
policy exclusion, support and a target compatible with their affected population.

An observed portfolio covariance matrix does not by itself identify endogenous
counterfactual portfolio choices. Even if a reduced-form shock loading $B_0$
is known perfectly, a new admissible-asset set may produce any loading $B_1$
allowed by unspecified preferences, pricing and constraints. The relevant
identified set must range over those structural primitives and equilibrium
selectors, not hold the baseline portfolio fixed by default.

The results here supply exact checks for an eventual three-region model: its
complete-market symmetric limit should reproduce (3)--(4); its CARA interior
limit should reproduce (5); its reported variance effect must state its
projection convention. Identifying safe-asset demand, exchange-rate valuation,
nominal adjustment and distributional welfare remains the unsolved part.
\begin{thebibliography}{9}
\bibitem{source} R. Kekre and M. Lenel (2024),
\emph{The Flight to Safety and International Risk Sharing}, AER 114(6), 1650--1691.
\url{https://www.aeaweb.org/articles?id=10.1257/aer.20211319}.
The publisher abstract describes safe-dollar demand, unequal risk bearing and
nominal rigidities. The finite-state complete-market benchmarks above are
separate calculations, not solutions of that monetary model.
\end{thebibliography}

\end{document}
